% ============================================================
%  Section 2: Related Work
%  AgentForge: A Multi-Agent Framework for Autonomous
%  Software Engineering
% ============================================================

\section{Related Work}
\label{sec:related}

\subsection{LLM-Based Code Generation}

Early work on neural code generation focused on
sequence-to-sequence models trained on code
corpora~\cite{iyer2016summarizing,yin2017syntactic}.
The introduction of large-scale pretraining on
mixed code and text data produced a step change in
capability: Codex~\cite{chen2021codex} demonstrated
that a GPT-style model trained on GitHub could solve
a substantial fraction of competitive programming
problems in a single forward pass. Subsequent
models — AlphaCode~\cite{li2022alphacode},
StarCoder~\cite{li2023starcoder}, and
CodeLlama~\cite{roziere2023codellama} — pushed
performance further through larger training sets,
improved tokenization, and fill-in-the-middle
objectives designed for code editing.

A critical limitation shared by all single-shot
approaches is their lack of mechanisms for
verification or iterative revision. When a generated
function fails a test case, the model has no channel
to observe the failure and produce a correction.
This limitation has motivated a paradigm shift from
single, monolithic LLMs to multi-agent systems,
where specialized agents handle planning, coding,
testing, and debugging. AgentForge embraces this
shift through its Tester--Debugger loop, which
provides genuine execution feedback at every step.

\subsection{Program Repair and Iterative Refinement}

Automated program repair (APR) is a classical
software engineering problem that predates deep
learning~\cite{genprog,angelix,prophet}.
Recent neural APR systems use LLMs to propose
patches given a failing test and a fault
localization signal~\cite{xia2022less,
jin2023inferfix,fan2023automated}. These systems
share with AgentForge the insight that execution
feedback is essential, but typically assume a
pre-existing test suite and a known fault location
rather than generating both from scratch.

Several recent systems have advanced iterative
repair with richer feedback signals.
TraceRepair~\cite{tracerepair2026} leverages
runtime execution traces as objective constraints
for patch validation, moving beyond simple error
messages. DynaFix~\cite{dynafix2025} introduces
execution-level dynamic information (variable
states, call stacks) to guide LLM patch generation,
achieving substantial improvements on Defects4J.
InspectCoder~\cite{inspectcoder2025} enables LLMs
to conduct dynamic analysis via interactive
debugger control, transforming debugging from
blind trial-and-error into systematic diagnosis.
RGD~\cite{rgd2024} employs three distinct LLM
agents (Guide, Debug, Feedback) to decompose the
code debugging process.

Self-repair~\cite{olausson2023self} and
self-debugging~\cite{chen2023teaching} frameworks
prompt an LLM to revise its own output given an
error message, without a separate specialized
agent. Our Debugger agent occupies a similar niche
but operates within a structured pipeline that
distinguishes code generation, test generation,
and repair as distinct subtasks with distinct
prompts — a separation we find important in our
ablation study (Section~\ref{sec:ablation}).

\subsection{Agentic and Tool-Using LLM Systems}

ReAct~\cite{yao2023react} introduced the paradigm
of interleaving reasoning traces (\textit{thoughts})
with tool invocations (\textit{actions}), enabling
LLMs to use web search, calculators, and code
interpreters within an open-ended loop.
Reflexion~\cite{shinn2023reflexion} extended this
with verbal reinforcement: the agent reflects on
its own failures and stores a natural language
summary in an episodic memory buffer.
Both systems operate with a single LLM instance
that must simultaneously plan, act, and evaluate.
AgentForge instead assigns these roles to distinct
agents with specialized prompts, following the
principle that specialization reduces the cognitive
load placed on any single model call.

Toolformer~\cite{schick2023toolformer} trains a
model to decide when and how to invoke external
tools by self-supervised annotation, removing the
need for explicit tool-use prompting.
While elegant, this approach requires fine-tuning
and does not provide the structured pipeline
guarantees that AgentForge offers for software
engineering tasks.

\subsection{Multi-Agent LLM Frameworks}

The idea of using multiple cooperating LLM instances
to solve complex tasks has attracted significant
recent attention.
MetaGPT~\cite{hong2023metagpt} encodes software
engineering roles — product manager, architect,
engineer, QA — as distinct agents that communicate
through structured documents, demonstrating that
role separation improves output quality on
software project benchmarks.
ChatDev~\cite{qian2023communicative} uses
chat-based communication between a CEO, CTO, and
programmer agent to produce complete small software
projects, including documentation and unit tests.
AutoGen~\cite{wu2023autogen} provides a general
framework for multi-agent conversation in which
agents can invoke tools, write and execute code,
and hand off to human proxies.

Recent frameworks have introduced self-evolution,
competitive collaboration, and structured
communication. SEMAG~\cite{semag2026} proposes
self-evolutionary agents that adapt their workflows
to task difficulty, achieving state-of-the-art
results on CodeContests. Eco-Evolve~\cite{ecoevolve2026}
introduces a dynamic agent topology with a dedicated
Critic agent and error-driven self-evolution via
Hindsight Experience Replay (HER), achieving strong
results on SWE-bench Verified. SWE-Debate~\cite{swedebate2026}
employs competitive multi-agent debate across three
rounds, leveraging fault propagation traces and
Monte Carlo Tree Search for patch generation.
COLLAB-LLM~\cite{collabllm2026} provides a
structured communication protocol and hierarchical
role architecture, enabling teams of up to eight
agents with reduced coordination overhead.

AgentForge differs from these systems in three
important ways.
First, it targets \textit{real-world bug fixing}
on actual GitHub repositories rather than synthetic
project generation or constrained benchmarks.
Second, it evaluates generated code in a
\textit{genuinely isolated} Docker sandbox with
strict resource caps and networking disabled,
providing authentic execution feedback that
cannot be fabricated — a mandatory step in
every pipeline run.
Third, it incorporates \textit{dual retrieval}:
episodic memory of past solved tasks (for
cross-task transfer) and a live repository index
(for intra-task grounding), a combination absent
from prior multi-agent frameworks.

\subsection{Autonomous Software Engineering Agents}

SWE-agent~\cite{yang2024sweagent} introduced an
agent--computer interface (ACI) that gives a
single LLM structured access to a shell, editor,
and search tools, achieving then-state-of-the-art
results on SWE-bench. Devin~\cite{cognition2024devin}
demonstrated a commercial system capable of
completing multi-day engineering tasks end-to-end,
though its architecture is not publicly disclosed.
OpenHands~\cite{wang2024openhands} (formerly
OpenDevin) provides an open-source platform
for building software engineering agents,
offering a sandboxed execution environment.
Trae Agent~\cite{trae2025} implements a test-time
scaling approach with generation, pruning, and
selection modules, attaining a Pass@1 score of
75.20\% on SWE-bench Verified. MAGIS~\cite{magis2024}
employs Manager, Repository Custodian, Developer,
and QA agents for GitHub issue resolution.

AgentForge is most directly comparable to
SWE-agent in scope and evaluation methodology,
but differs in its \textit{structured multi-agent
decomposition}: where SWE-agent uses a single
LLM interacting with shell tools, AgentForge
uses five specialized agents each with a
distinct role, a dedicated prompt, and a
clearly scoped output contract. Our ablation
study quantifies the value of each agent
independently, providing interpretable evidence
for the contribution of each design decision.

\subsection{Benchmarks for Software Engineering}

HumanEval~\cite{chen2021codex} and
MBPP~\cite{austin2021program} measure function
synthesis from docstrings — tasks that are well-suited
to single-shot generation but do not capture the
complexity of real-world bug fixing.
SWE-bench~\cite{jimenez2024swebench} addresses
this gap with 2,294 issues drawn from 12 popular
Python repositories, each paired with a gold
patch and a test suite that verifies the fix.
SWE-bench Lite reduces this to 300 carefully
filtered instances, balancing evaluation cost
with statistical power. SWE-bench Verified
provides a human-validated subset for more
reliable evaluation. Defects4J~\cite{defects4j}
is a widely used benchmark for Java program
repair, employed by systems such as TraceRepair.
We evaluate on SWE-bench Lite as our primary
benchmark, following recent convention in the
agentic software engineering literature~\cite{yang2024sweagent,
wang2024openhands,zhang2024autocoderover}.

\subsection{Memory and Retrieval in LLM Systems}

Retrieval-augmented generation
(RAG)~\cite{lewis2020rag} demonstrated that
augmenting LLM inputs with retrieved passages
substantially improves factual accuracy on
knowledge-intensive tasks. In the code domain,
repository-level code completion
systems~\cite{zhang2023repocoder,shrivastava2023repofusion}
retrieve relevant functions and files to provide
context for the generation model.
GraphCodeAgent~\cite{graphcodeagent2025} introduces
a dual-graph approach (Requirement Graph and
Structural-Semantic Code Graph) to guide multi-hop
retrieval for repository-level tasks.

AgentForge extends this paradigm with a
\textit{dual retrieval} strategy: an episodic
memory of past solved tasks (for cross-task
transfer) and a live repository index (for
intra-task grounding), both backed by the same
ChromaDB vector store and queried by cosine
similarity of OpenAI text embeddings.
To our knowledge, this is the first system to
combine both retrieval sources within an
agentic software engineering pipeline.